\chapter{Experiments \& Statistical Inference}\label{ch:experiments-and-inference}

% Scope: null-hypothesis testing, statistical power & sample size, OLS regression,
%   A/B-test design for business decisions. NEW chapter (closes the experiment-design gap).
% The book uses A/B tests and incrementality RCTs (\cref{ch:attribution-and-measurement})
%   without teaching the machinery; this chapter supplies it, extending \cref{ch:decisions-under-uncertainty}.
% Draft per house style: flowing prose + example/admonition set-offs; running-SaaS worked examples.

\section{The Null-Hypothesis Framework}\label{sec:null-hypothesis}
% complexity: 2

Every product decision is a bet. The null-hypothesis framework forces the decision-maker to
quantify how unlikely the observed data would be if the intervention had \emph{no} effect,
then decide whether that improbability is low enough to act. The managerial intuition is
simple: collect evidence, ask how often data this extreme would appear by chance under the
status quo, and reject the status quo only when that chance falls below a pre-committed
threshold. The discipline is in the pre-commitment: choosing the threshold before seeing
results, not after.

Under the null hypothesis \(H_0: \mu_1 = \mu_2\) that two groups have the same mean, the
test statistic for large samples is a standardised distance between the observed group means,
scaled by the standard error of their difference:
\begin{equation}\label{eq:z-test}
  z \;=\; \frac{\bar{x}_1 - \bar{x}_2}%
               {\sqrt{\dfrac{\sigma_1^2}{n_1}+\dfrac{\sigma_2^2}{n_2}}}.
\end{equation}
Reject \(H_0\) when \(|z| > z_{\alpha/2}\), the critical value for a two-tailed test at
significance level \(\alpha\). For the conventional \(\alpha=0.05\), \(z_{\alpha/2}=1.96\).
For proportion outcomes---conversion rates, trial-to-paid rates---\(\sigma_i^2 = p_i(1-p_i)\).
The formula rests on the Central Limit Theorem, which requires samples large enough for the
sampling distribution of the mean to be approximately normal; in practice, \(n \geq 30\)
per group is a common floor, though for small proportions much larger samples are needed
before the approximation is reliable.

\begin{example}[SaaS onboarding A/B test]
The product team tests a redesigned onboarding flow. Control: \num{200} trials, \num{50}
convert to paid (\(\hat{p}_c = \qty{25}{\percent}\)). Treatment: \num{200} trials, \num{62}
convert (\(\hat{p}_t = \qty{31}{\percent}\)). The decision: is the \qty{6}{\percent} lift
real or noise?

Set \(\alpha=0.05\), so the critical value is \(z_{0.025}=1.96\). Estimate the pooled
proportion under \(H_0\) as \(\hat{p} = (50+62)/(200+200) = 0.28\), giving pooled
\(\hat{\sigma}^2 = 0.28 \times 0.72 = 0.2016\). The standard error is
\[
  \mathrm{SE} \;=\; \sqrt{\frac{0.2016}{200}+\frac{0.2016}{200}}
               \;=\; \sqrt{0.002016}
               \;\approx\; 0.0449.
\]
Applying \cref{eq:z-test},
\[
  z \;=\; \frac{0.31 - 0.25}{0.0449} \;\approx\; 1.34.
\]
Since \(|z|=1.34 < 1.96\), the result is \emph{not} statistically significant at \(n=200\)
per arm. The \qty{6}{\percent} lift could easily be chance. The implication is not that the
treatment fails; it is that the test is too small to distinguish a real effect of this
magnitude from sampling noise. Collecting more data is the only remedy---how much more is
the subject of \cref{sec:power-and-sample-size}.
\end{example}

\begin{admonition}{Warning}
The p-value is the probability of seeing a test statistic at least this extreme \emph{if the
null were true}. It is emphatically not the probability that the null is true, nor that the
alternative is false. A p-value of \num{0.03} does not mean there is a \qty{97}{\percent}
chance the treatment works; it means the data are moderately incompatible with ``no effect.''
Equally, statistical significance says nothing about effect size: a trivial \qty{0.1}{\percent}
lift will reach \(p<0.05\) with a large enough sample. Business decisions require both a
significant result \emph{and} an effect large enough to matter economically.
\end{admonition}

The test statistic in \cref{eq:z-test} provides the machinery; \cref{sec:power-and-sample-size}
establishes how many observations are needed before the machinery is powerful enough to be
worth running. \textcite{anderson2019statistics} develop the full sampling-distribution
theory; \textcite{wooldridge2020introductory} extends it to regression settings where
multiple covariates enter simultaneously.


\section{Statistical Power \& Sample Size}\label{sec:power-and-sample-size}
% complexity: 4

An experiment that cannot detect the effect it was designed to detect is not a neutral
non-result; it is a waste. Statistical power---the probability of correctly rejecting a false
null---quantifies the risk of that failure. Low power means the experiment will most often
return an inconclusive result even when the treatment genuinely works, producing a stream of
``no significant difference'' findings that are misread as evidence the product is fine. The
business cost is symmetric: underpowered tests waste time on real improvements and give false
reassurance when a change is causing harm.

Power depends on three inputs the experimenter controls: the sample size \(n\), the
significance level \(\alpha\), and the minimum detectable effect \(\delta\). For a two-arm
test with equal group sizes and common variance \(\sigma^2\), the required sample per arm is:
\begin{equation}\label{eq:sample-size}
  n \;=\; \frac{2\,(z_{\alpha/2}+z_{\beta})^{2}\,\sigma^{2}}{\delta^{2}}.
\end{equation}
Here \(z_{\beta} = \Phi^{-1}(1-\beta)\) is the critical value for the desired power
\(1-\beta\); at \qty{80}{\percent} power, \(z_{\beta}=0.84\). The inverse-square relationship
between \(\delta\) and \(n\) is the central fact of experimental design: halving the effect
size you care about \emph{quadruples} the required sample size. The intuition is that smaller
signals are harder to distinguish from noise, so you need more observations to drown out
that noise. \Cref{fig:power-curve} plots power against effect size for three sample
sizes, making the trade-off concrete.

\begin{figure}[htbp]
  \centering
  \includegraphics[width=0.86\linewidth]{power-curve}
  \caption{Statistical power versus absolute effect size (conversion-rate lift over a
    \qty{25}{\percent} baseline) at \(\alpha=0.05\), two-tailed, for \num{100}, \num{200},
    and \num{550} observations per arm. The dashed line marks the \qty{80}{\percent} power
    target; the diamond marks the required design point. Detecting the \qty{6}{pp} lift from
    \cref{sec:null-hypothesis} at \qty{80}{\percent} power requires roughly
    \num{879} observations per arm.}
  \label{fig:power-curve}
\end{figure}

\begin{example}[Sizing the onboarding A/B test]
How many trials are needed to detect the \qty{25}{\percent} \(\to\) \qty{31}{\percent}
onboarding lift at \qty{80}{\percent} power, \(\alpha=0.05\)? The minimum detectable effect
is \(\delta=0.06\). Use the pooled proportion \(\hat{p}=0.28\) to estimate
\(\sigma^2 = 0.28\times 0.72 = 0.2016\), and plug into \cref{eq:sample-size}:
\[
  n \;=\; \frac{2\,(1.96+0.84)^{2}\times 0.2016}{0.06^{2}}
      \;=\; \frac{2\times 7.84 \times 0.2016}{0.0036}
      \;\approx\; \num{879} \text{ per arm}.
\]
The funnel delivers \num{800} trials per month, so with \num{400} per arm the experiment
reaches \num{879} after roughly \num{2.2} months---a manageable horizon. Is it worth the
wait? The value of confirming the lift: \qty{6}{pp} more paid conversions \(\times\) \num{200}
new customers/month \(\times \money{3150}\) CLV per customer \(= \money{37800}\) of incremental
monthly contribution---a recurring stream worth \(\money{37800}/0.08 \approx \money{472500}\)
in perpetuity at the firm's \qty{11}{\percent} WACC less \qty{3}{\percent} growth
(\cref{eq:gordon}). Against a two-month delay, that ratio overwhelmingly favours running a
disciplined test over guessing and potentially shipping a change that does nothing or causes
harm.
\end{example}

\begin{admonition}{Warning}
Optional stopping---peeking at results daily and stopping the moment \(p<0.05\)---inflates
the true Type-I error far above the nominal \qty{5}{\percent}. If you peek ten times during
a trial, the probability of a spurious significant result can exceed \qty{30}{\percent}. The
fix is straightforward: pre-register the stopping rule and the required \(n\) before launch,
and read results exactly once. Platforms that offer ``always-valid'' sequential testing
(\textcite{johari2022always}) allow continuous monitoring without inflating error, but they
require a different underlying test statistic (a mixture sequential probability ratio test,
mSPRT) and typically need larger samples than the fixed-horizon formula above to achieve
the same power.
\end{admonition}

The sample-size calculation in \cref{eq:sample-size} is the pre-commitment that makes
\cref{eq:z-test} honest. Without it, a team can always find significance by running the
experiment longer---a practice that destroys the statistical validity of everything it
touches.


\section{Regression: Controlling for Covariates}\label{sec:regression}
% complexity: 3

Randomised experiments settle causation by construction: random assignment balances all
observed and unobserved confounders between treatment and control, so any difference in
outcomes must reflect the treatment. But many business questions cannot be randomised:
churn depends on tenure, product depth, plan tier, and sales channel, and customers were not
randomly assigned to those states. Ordinary least squares (OLS) regression answers the
question ``how much does each factor matter, holding all others fixed?'' It is the
workhorse of observational business analysis precisely because it is the simplest tool for
holding some things constant while examining others.

The multiple linear regression model posits that the outcome \(y_i\) is a linear function
of \(k\) explanatory variables plus an error term capturing everything the model omits:
\begin{equation}\label{eq:ols}
  y_i \;=\; \beta_0 + \beta_1 x_{1i} + \cdots + \beta_k x_{ki} + \varepsilon_i,
  \qquad
  \hat{\boldsymbol{\beta}} \;=\; (\mathbf{X}^{\!\top}\mathbf{X})^{-1}\mathbf{X}^{\!\top}\mathbf{y}.
\end{equation}
The matrix solution \(\hat{\boldsymbol{\beta}}\) minimises the sum of squared residuals. By
the Gauss--Markov theorem, OLS is the Best Linear Unbiased Estimator (BLUE) when the errors
have zero mean, constant variance (homoskedasticity), and are uncorrelated with the
regressors. Each coefficient \(\hat{\beta}_j\) measures the marginal change in \(y\) for a
one-unit change in \(x_j\), \emph{ceteris paribus}. That ``all-else-equal'' qualifier is the
source of both the method's power and its danger: it holds constant only what is in the
model, and unobserved confounders absorbed into \(\varepsilon\) will bias the estimates if
they correlate with the regressors.

\begin{example}[Quantifying the integration switching-cost moat]
The customer-success team wants to know how much product depth protects against churn. They
regress monthly churn rate on three customer-level variables: tenure in months, number of
active integrations, and plan tier (1~=~starter, 2~=~professional, 3~=~enterprise). The
estimated model on the firm's \num{1200}-customer cohort returns:
\begin{align*}
  \widehat{\text{churn}}_i \;=\;{}
    & 0.032
      - 0.00012\,\text{tenure}_i
      - 0.0080\,\text{integrations}_i
      - 0.0031\,\text{tier}_i.
\end{align*}
The coefficient on integrations is \(\hat{\beta}_{\text{int}} = -0.008\): each additional
active integration reduces monthly churn by \qty{0.8}{pp}, holding tenure and tier fixed.
At the firm's baseline monthly churn of \qty{0.65}{\percent}, a customer with five
integrations instead of one has predicted churn of
\(0.65\% - 4\times0.8\% = 0.65\% - 3.2\% \approx \) effectively zero---a near-perfect
retention profile. This coefficient quantifies the switching-cost moat described in
\cref{sec:moat-taxonomy}: every integration the customer builds is a thread binding them
to the platform. The implication for product investment: a feature that encourages customers
to add one more integration is worth, in expected annual CLV, the reduction in churn times
the per-customer CLV of \money{3150}.
\end{example}

\begin{admonition}{Warning}
Regression coefficients are correlational, not causal, unless the assignment mechanism is
randomised or a valid instrument is used. The negative correlation between integrations and
churn could partly reflect \emph{selection}: customers who build deep integrations may have
chosen the product because they planned to stay, not merely because integrations raise
switching costs. The only way to isolate the causal contribution of integrations is a
randomised experiment---nudging a random subset of users toward additional integrations and
comparing subsequent churn (\cref{sec:ab-test-design}). Observational regression controls
for measured confounders; randomisation controls for everything.
\end{admonition}

\textcite{wooldridge2020introductory} prove the Gauss--Markov result and detail the
diagnostics needed to assess whether its assumptions hold; \textcite{angrist2009mostly}
connect OLS to causal inference and explain when instrumental variables or difference-in-
differences designs are needed instead. The switching-cost coefficient above is one example
of a broader class of value-creation mechanisms catalogued in \cref{sec:moat-taxonomy};
the same regression framework applied to geo-level data underlies the incrementality
measurement in \cref{sec:incrementality}.


\section{A/B-Test Design for Business Decisions}\label{sec:ab-test-design}
% complexity: 3

Knowing the mechanics of \cref{eq:z-test} and \cref{eq:sample-size} is necessary but not
sufficient. A well-designed test requires four upfront decisions that have nothing to do with
the formula: what primary metric to optimise, which guardrail metrics must not regress, at
what level to randomise users, and how long to run. Skipping any of these turns a statistically
valid calculation into a commercially unreliable experiment. The practitioner's discipline is
to write down all four before any data are collected.

The \textbf{primary metric} (overall evaluation criterion, OEC) should be the metric most
directly tethered to long-run value. For a SaaS onboarding test, trial-to-paid conversion
is the right primary metric: it directly determines the \num{200} new customers per month in
the fact sheet and, through CLV, the \money{37800}/month lift computed in
\cref{sec:power-and-sample-size}. Revenue per trial is tempting but too noisy; session-to-trial
is upstream and may move without downstream benefit. \textbf{Guardrail metrics} protect core
business health: if the treatment improves trial-to-paid but slows page load by two seconds,
the short-term gain will be erased by reduced long-run engagement. Pre-specifying guardrails
prevents the team from launching a treatment that wins on the primary metric while silently
damaging the funnel.

\textbf{Randomisation unit} determines which guardrail violations are detectable and governs
whether the Stable Unit Treatment Value Assumption (SUTVA) holds. Randomising at the session
level is straightforward but can contaminate results if the same user lands in both arms
across sessions. For an onboarding test, user-level or device-level assignment is the correct
choice: the treatment is the full onboarding experience, not a single page view, and the
conversion event occurs across multiple sessions. Session-level randomisation would mix
control and treatment experiences for the same trial, biasing the observed conversion rate
in both arms. Where network effects mean that one user's outcome depends on another user's
treatment status---common in social products and two-sided marketplaces---geo-level
randomisation (as in the incrementality framework of \cref{sec:incrementality} and the iROAS
estimator in \cref{eq:iroas}) is the only design that avoids SUTVA violations.

\begin{example}[Designing the pricing-page test]
The growth team wants to test a pricing-page redesign that makes the professional plan more
prominent. The decision framework is:

\begin{description}
  \item[Primary metric] Trial-to-paid conversion rate. Current baseline \qty{25}{\percent};
    minimum detectable effect \qty{4}{pp} (a lift from \qty{25}{\percent} to \qty{29}{\percent}),
    anchoring to the expected margin from \num{8} incremental paid conversions per month
    (\(0.04\times 200\)) times \money{3150} CLV \(= \money{25200}\)/month.
  \item[Guardrail metrics] (a) Session-to-trial rate: must not fall by more than \qty{2}{pp}
    (a \qty{10}{\percent} relative drop would eliminate the funnel gain); (b) Support-ticket
    volume: no more than a \qty{15}{\percent} increase (a heavily interrogated pricing page
    inflates cost-to-serve).
  \item[Randomisation unit] User (cookie + login), assigned on first session. A user in the
    treatment arm sees the redesigned pricing page on every visit; a control user always sees
    the original.
  \item[Sample size and duration] Using \cref{eq:sample-size} with \(\delta=0.04\),
    \(\sigma^2=0.25\times 0.75=0.1875\), \(\alpha=0.05\), \(1-\beta=0.80\):
    \[
      n \;=\; \frac{2\,(1.96+0.84)^{2}\times 0.1875}{0.04^{2}} \approx \num{1838}
        \text{ per arm.}
    \]
    At \num{800} trials per month split equally across two arms (from the funnel in the
    fact sheet), the test requires roughly \num{4.6} months of data---a meaningful wait. The team must either accept a larger MDE
    (say, \qty{6}{pp} \(\to\) \num{879} per arm \(\approx\) 2.2 months) or invest in reducing
    conversion variance through CUPED-style pre-experiment covariate adjustment. Pre-specifying
    the stopping rule to three weeks at \num{879}-per-arm and reading results exactly once
    eliminates optional-stopping bias.
\end{description}

The implication: the binding constraint on experimentation velocity is not statistical theory
but funnel size. A firm that acquires only \num{200} paid customers per month needs months to
run even modestly powered tests on downstream conversion. Investing in top-of-funnel
acquisition (\cref{ch:decisions-under-uncertainty}) expands the experimental surface area as
much as it expands revenue.
\end{example}

\begin{admonition}{Note}
Pre-registering the design---primary metric, guardrails, randomisation unit, sample size,
and stopping rule---before launch is not bureaucratic overhead; it is the mechanism that
keeps the p-value from the analysis honest. Without pre-registration, a team can always
find a metric that moved, a time window that shows significance, or a subgroup that
confirms their prior. \textcite{kohavi2020trustworthy} document the most common failure
modes of large-scale A/B programmes at Microsoft and elsewhere; \textcite{gordon2019comparison}
show that observational methods can miss true causal effects by large, inconsistent margins,
which is precisely why a properly designed test---even a slow one---is worth running when
the revenue at stake justifies the iROAS calculation in \cref{eq:iroas}.
\end{admonition}

The test design framework here is the operational complement to the incrementality toolbox
in \cref{sec:incrementality}: both ask the same causal question---does this intervention
produce more revenue than the counterfactual?---but at different granularities. Product
experiments randomise users within a product surface and measure conversion; geo-holdout
experiments randomise geographic markets and measure the full-funnel iROAS of media spend.
\textcite{johnson2017ghost} demonstrate the magnitude of advertiser overstatement when
geo-level randomisation is bypassed in favour of self-reported platform attribution. The
connecting thread between all three chapters---this chapter, \cref{ch:attribution-and-measurement},
and \cref{ch:decisions-under-uncertainty}---is the potential outcomes framework: a causal
effect exists only as a comparison against a counterfactual that was never observed, and the
only reliable way to construct that counterfactual is randomisation.
